\songtitle{Luciafestens dugg}

\tag*{2026}\tag{thomas-entourage}\tag{erika-czajka}\tag{thomas-garcia}\tag{swedish-lyrics}
\begin{notes}
Adaptation into Swedish by Claude of \emph{That Night}.
\end{notes}
\begin{style}
Vispop with lilting nyckelharpa and bright fiol over a midtempo pulse, buoyant but restrained, male baritone lead vocals, close-mic delivery, swedish lyrics
\end{style}
\begin{lyrics}
\annotation{Lyrics}

\songpart{Verse 1}
Du var Niklas flickvän, det brydde mig knappt,\\
ditt lätta sätt, din blick, jag bara stannade kvar.\\
Du gled på isen medan jag höll mig tappt,\\
utan att håna hur klumpig jag var.

\songpart{Pre-Chorus}
De djupgrå ögonen bakom din lugg,\\
Luciafesten i Stockholm, ditt leende där.\\
Jag lät mig försvinna i tankarnas dugg,\\
du visste det ändå, och log som du är.

\songpart{Chorus}
Vi pratade den natten, så lätt, så rätt,\\
genom gränser, genom avstånd,\\
samtalet tog aldrig slut,\\
aldrig slut, åh Erika, aldrig slut.

\songpart{Verse 2}
Sportlov i Tallinn, du sa att du kom,\\
du och Niklas var historia, som om det spelade roll.\\
Vi pratade om städer, om ett grönare rum,\\
att bygga för folk, hålla dem på samma håll.

\songpart{Pre-Chorus}
Och pratet blev en kram, och kramen blev en kyss,\\
ett tecken att våra ord alltid skulle länka.\\
Din kropp, min kropp, omöjliga att missa,\\
en gräns vi korsade innan vi hann tänka.

\songpart{Chorus}
Vi älskade den natten, så lätt, så rätt,\\
genom gränser, genom avstånd,\\
samtalet tog aldrig slut,\\
aldrig slut, åh Erika, aldrig slut.

\songpart{Bridge}
Jag kände ditt sinne, jag kände din glöd,\\
det samtalet i Tallinn var ditt hela cv.\\
Det där är ingen plats för en älskare —\\
snarare en plats för elden inom dig,\\
en eld jag känner nära, min vän,\\
smidd på is, och kärleken till det gröna.

\songpart{Chorus}
Vi kände den natten, så lätt, så rätt,\\
genom gränser, genom avstånd,\\
samtalet tog aldrig slut,\\
aldrig slut, åh Erika, aldrig slut.

Det slutade inte, min vän, aldrig slutade.
\end{lyrics}

